\section{威胁模型：攻击者能力、控制面与后果}

攻击面告诉我们哪里可能失守，威胁模型则规定攻击者实际能触碰哪些位置、投入多少资源以及希望造成什么后果。两者不能颠倒：如果先选一个热门 benchmark 再补 threat model，评测可能只证明系统能抵抗弱而不现实的攻击。这里应把 attacker capability、knowledge、persistence、access channel 和 consequence 写成明确假设，再决定测试强度与防线。

讲者强调了 Threat Modeling 的必要性：不同攻击者能力下，最优防御完全不同。课程给出可操作框架：先定义攻击者可见面与控制面，再定义攻击目标和攻击后果，最后映射到检测与阻断点。

\begin{table}[H]
\centering
\begin{tabular}{p{2.2cm}p{4.1cm}p{4.6cm}}
\toprule
\textbf{威胁模型} & \textbf{攻击者能力} & \textbf{常见后果} \\
\midrule
Black-box & 仅可通过公开接口输入内容 & 注入、越权诱导、间接命令执行 \\
Gray-box & 知道部分系统结构与工具接口 & 构造高命中攻击链、规避浅层护栏 \\
White-box & 可读代码/配置/提示词 & 精准定位注入点、自动化生成攻击样本 \\
\bottomrule
\end{tabular}
\caption{三类常见威胁模型与攻击收益}
\end{table}

\begin{importantbox}{课程中的核心判断}
如果没有清晰 Threat Model，评测指标会失真，防御策略会错位。很多 ``防住了'' 的结果，只是在弱攻击者假设下成立。
\end{importantbox}

Black-box 并不等于低风险。攻击者可以长时间自适应查询、创建多个账号、控制网页或邮件内容，并观察 Agent 的外部行动；这些反馈足以支持 fuzzing。Gray-box 攻击者可能从公开文档、错误信息、tool schema 或前员工知识推断内部结构。White-box 则可以直接搜索 prompt、policy 和代码路径。评测报告必须写明查询预算、是否可观察工具结果、是否允许污染持久状态、攻击是否跨会话，以及攻击者是否知道防御策略。

后果也要按实际损失分级。一次 system prompt 泄露可能只是信息暴露，也可能揭示内部工具和绕过条件；一次错误网页点击可能无害，也可能下载恶意文件；一次越权数据库读取可能只触及测试记录，也可能横向访问整个客户租户。应以 confidentiality、integrity、availability、financial loss、physical harm 和 recovery cost 共同描述 consequence，而不是用单一 ASR 把所有成功攻击视为同等严重。

\begin{warningbox}{威胁模型必须包含组合与时间}
真实攻击者会先侦察、再污染、等待 memory 持久化，最后在高权限任务中触发 payload。只测试单轮 prompt 会漏掉跨会话、跨工具和延迟触发攻击；只测试最强攻击又可能掩盖低成本批量滥用。至少应覆盖 opportunistic、targeted 和 persistent 三档对手。
\end{warningbox}

\subsection{本章小结}
威胁建模不是文档工作，而是安全工程入口。先定义攻击者，再定义防御；先定义后果，再定义指标。

